The current project deliberately favors scope control over breadth.
That is a strength for the first paper draft, but it also creates limitations
that should be stated clearly.

\paragraph{Current limitations.}
\begin{itemize}[leftmargin=1.25em]
  \item The active implementation currently covers Stage 1 and Stage 2 only.
  \item The task uses a single provided evidence passage rather than retrieved
  evidence.
  \item The action space is intentionally narrow and does not yet cover richer
  support labels or multi-hop settings.
  \item The current results will reflect one controlled dataset setup rather
  than a broad claim about hallucination in general.
\end{itemize}

\paragraph{Intended deployment scope.}
The current approach is best interpreted as a method for high-trust,
document-bounded QA, not as a general-purpose opinion or advice generator.
Its most plausible deployment settings are ones in which the evidence source is
explicit, the answer should remain short and verifiable, and abstaining is
preferable to guessing.
If the provided evidence is incomplete, outdated, or systematically biased, the
method may still produce a well-grounded but incomplete answer, so groundedness
should not be confused with neutrality or full real-world balance.

\paragraph{How to interpret the current result.}
The current evidence supports two related but different claims.
First, the Stage 1 run shows that explicit abstention supervision and
abstention-aware evaluation can strongly reduce unsupported answers in a simple
extractive QA system.
Second, the Stage 2 run shows that support verification is learnable on top of
that abstain-aware base.
However, the current verifier mostly predicts support rather than intervening
aggressively enough to shift end-to-end QA behavior.
The result should therefore be read as a trade-off improvement plus a partially
successful extension, not yet as a generally better grounded QA model.

\paragraph{Why this stage is still useful.}
Even under its narrow scope, the first two stages now test two defensible
claims with clean supervision and interpretable failure modes.
Stage 1 shows that abstention can be taught.
Stage 2 shows that support can be predicted, but also reveals the remaining
bottleneck clearly:
the system still needs better calibration and tighter decision control before a
support signal becomes a strong downstream filter.

\paragraph{Legacy plan and revision.}
The original roadmap after Stage 1 was to move directly from Stage 2 support
verification into Stage 4 unsupported-confidence control, while treating
Stage 3 calibration as supporting work and deferring retrieval to a later
realism check.
That legacy plan assumed that once support became learnable, the raw verifier
scores would already be good enough to drive downstream control.
Stage 2 changed that assumption.
On the dev split, support recall reaches 99.43\%, but support precision is only
66.62\%, the verifier predicts support for 98.88\% of answered cases, and the
gated system improves overall \(F_1\) by only 0.03 points.
So the long-term goal stays the same, but the immediate plan changes:
calibration now becomes the next dependency rather than optional supporting
work, because the current support scores are still too permissive for a strong
fixed controller.

\paragraph{Next paper-facing extensions.}
With Stage 1 and Stage 2 now in place, the most natural paper extensions are:
\begin{itemize}[leftmargin=1.25em]
  \item calibrating confidence against correctness and support,
  \item making the verifier less permissive on unsupported answers,
  \item testing retrieval noise, and
  \item comparing adaptive balancing against simpler fixed baselines.
\end{itemize}
